% !TeX root = ../main.tex
%%% ---------------------------------------------------------------------------
%%% 结论
%%%
%%% 结论不是摘要的复述。GB/T 7713.2 要求结论「明确、精练、完整、准确」，
%%% 应当包含：
%%%   1. 本研究得到的结果说明了什么（要有具体数据支撑）
%%%   2. 与已有结论的异同
%%%   3. 本研究的局限
%%%   4. 尚待解决的问题与后续方向
%%% 不要在结论里引入新的实验数据或新的论点。
%%% ---------------------------------------------------------------------------
\section{结论}
\label{sec:conclusion}

针对遥感图像中小目标特征在深层网络中衰减的问题，本文提出了跨层注意力
门控与可变形对齐相结合的特征融合方法。理论分析表明，门控融合在特征
互不相关的假设下不会放大方差（\cref{thm:variance}），因而不损害训练
稳定性；\dataset{} 上的实验表明，所提方法将 mAP 从 \qty{74.7}{\percent}
提升至 \qty{78.4}{\percent}，小目标子集的平均精度提升
\qty{6.2}{\percent}，推理速度由 \qty{33}{\fps} 降至 \qty{31}{\fps}，
代价可以接受。

本研究存在两点局限。其一，\cref{assum:indep} 的独立性假设在相邻层级
特征高度相关时并不成立，此时\cref{eq:var-bound} 的界会变松；其二，
实验仅在光学遥感图像上验证，对 SAR 等其他成像模态的适用性尚未考察。

后续工作将从两方面展开：一是把门控权重的生成从通道统计量扩展到显式的
尺度先验，二是在多模态遥感数据上验证方法的迁移能力。
